\section{Pruning and inlining the arithmetic network}
\label{sec:height23}
The new height-one improvements change the normalized arithmetic DAG before
network compilation. They do not require a new incidence-energy identity.
The graph construction, equality penalties, unit anchoring and universal
lift~\eqref{eq:networklift} continue to apply to the changed normalization.
The original four arithmetic source families remain those identified by
\file{height_one/SOURCE_BINDINGS.json}; a lower direct-quartic dimension
elsewhere is not silently substituted for one of these source theorems.

\subsection{Fuse source products and prune unused constructors}
Start with the retained triangular monic constant and scale definitions.
For a constant source occurrence of value $c$, replace $K_c$ by $K_aK_b$
when $ab=c$ and both factors are available. For a linear occurrence $cx$,
use an available $S(c,x)$ or $K_aS(b,x)$ with $ab=c$. Signs are retained.
The inherited factorization of quadratic source monomials is unchanged.
Every replacement is one signed monomial of degree at most two and has the
same exact pullback to the original source coordinates.

These are identities after constructor substitution, not arbitrary-tuple
identities: $K_6$ and $K_2K_3$ can differ away from constructor zeros.
Retain only constructors reachable from the newly selected source terms,
using the actual definition dependencies, and compact in inherited order.
Every retained dependency is then retained, so monic triangular induction
still uniquely evaluates all new constructors over every original tuple.
Every deleted old constructor has a unique intended value as well.

\begin{proposition}\label{prop:network-prune}
The replacement and pruning preserve the normalized integer zero set by
inverse projection and canonical restoration maps. The corresponding final
network zero sets are bijective through their common original source.
\end{proposition}
\begin{proof}
A zero of the old normalization projects to its original tuple, on which
every replacement has the same pullback. The retained constructors have
their intended values, so its projection is a new zero. Conversely, at a
new zero monic induction restores the original source equations and all
retained constructor values. Evaluate every deleted constructor at its
unique intended value; then every old row vanishes. These maps are inverse
on zero sets. Compose them with the already proved unique network lifts and
projections. No sign restriction, division or congruence condition is added.
\end{proof}

The finite selector uses product-first ordering, then at most five ordered
strict improvements of the number of reachable constructors and their total
occurrence count. Those costs are heuristics, not the final support. The
actual final collection determines the measured improvement.

\subsection{Inline low use constants by addition}
Let $K_c$ be retained, and choose $a+b=c$ with $1\le a\le b<c$.
Use $K_1=1$ as notation, not a new coordinate. Other factors are earlier
retained constants. Delete the defining row of $K_c$, substitute $K_a+K_b$
in every other row, and expand distributively.

Every replacement remains degree at most two. If a quadratic monomial
contains $K_c$ twice, retain all four signed occurrences. In particular
$a=b$ creates repeated occurrences, not a collected coefficient-one claim.
Every remaining constructor keeps its own coefficient-one pivot, and its
new dependencies are earlier, so triangular monic induction still gives its
original intended value. Each source pullback is unchanged since $a+b=c$.
A new zero uniquely restores $K_c=c$; all old rows and its deleted definition
then vanish. Deletion and restoration are inverse over both $\R$ and $\Z$.
Repeated inlining composes the restorations in reverse deletion order.
The usage budget and pair order choose valid substitutions but do not prove
monotone final support or optimality.

\subsection{Expand rare scale coordinates}
For a retained scale $S(c,x)$, $c>1$, delete its monic defining row and
replace every occurrence by the linear polynomial $cx$, represented as
$c$ signed copies of $x$. A quadratic occurrence containing the scale
twice gives $c^2$ equal signed degree-two occurrences. All source coordinates
precede constructors, so the other pivots stay monic and triangular.
Their exact pullbacks are unchanged. Restore the deleted scale uniquely as
$cx$, for negative, zero and positive $x$ alike. The same inverse zero-set
maps therefore apply. Prune newly unused constructors by
Proposition~\ref{prop:network-prune}.

The finite rule tests the precise increase of normalized occurrences against
a budget. The subsequent graph may have more or fewer monomials; neither
direction is a theorem. Repeated occurrences are represented as distinct
bound source occurrences, with factor aliases chosen as needed. They are
not confused with final collected coefficients. Height one is required and
checked only for the entire collected polynomial.

\subsection{The new finite catalog}
Every entry below is a polynomial on unrestricted integers, has exact degree
four, constant one, coefficient set $\{-1,+1\}$, and is nonnegative on all
real tuples. There are no equality dummies and every stated coordinate is
live. Support includes the constant. The source numbers are inherited family
identifiers, not a claim that the normalized source has that dimension.
\begin{center}\small
\begin{tabular}{@{}lrrrrr@{}}\toprule
Endpoint & Variables & Terms & Aliases & Flows & Source\\\midrule
minimum484 & 484 & 24131 & 461 & 23 & 241\\
compact485 & 485 & 23698 & 462 & 23 & 241\\
compact486 & 486 & 20007 & 462 & 24 & 242\\
compact487 & 487 & 19748 & 463 & 24 & 242\\
compact488 & 488 & 19649 & 464 & 24 & 242\\
compact489 & 489 & 17729 & 462 & 27 & 242\\
compact493 & 493 & 17058 & 465 & 28 & 242\\
compact494 & 494 & 16557 & 464 & 30 & 242\\
compact495 & 495 & 16400 & 465 & 30 & 242\\
balanced531 & 531 & 12020 & 484 & 47 & 263\\
balanced533 & 533 & 11985 & 486 & 47 & 263\\
sparse599 & 599 & 8279 & 504 & 95 & 285\\
sparse600 & 600 & 8241 & 505 & 95 & 285\\
sparser695 & 695 & 7135 & 504 & 191 & 285\\
sparser696 & 696 & 7113 & 505 & 191 & 285\\
\bottomrule\end{tabular}
\end{center}


The new $484/24131$ variable/support point improves both dimensions against
the Round22 minimum $493/24828$. The new $533/11985$ point improves the
retained $557/11991$ tradeoff in both measures. Likewise $600/8241$ improves
$605/8305$, and $696/7113$ improves $701/7141$. Every one of the nine selected
Round22 endpoints is dominated by at least one of the selected Round23
endpoints in variable count and support, with degree and height fixed.
At exactly 493 variables, support falls from the old 24828 to 17058,
a reduction of 7770 monomials (about 31.3\%). This is separate from the
new minimum of 484 variables. The new $495/16400$ point also improves the
old $496/20548$ point.
This is a checked finite comparison, not an exhaustive global frontier.
Within the new catalog the alternatives trade dimension for support.

\subsection{Independent source semantics and final collection}
The independent normalizer audit computes each coordinate's intended
polynomial from metadata without importing the producer. It verifies every
source row, with its certified sign, has exactly its pinned pullback;
every constructor has identically zero pullback and exactly one monic
triangular pivot; and every source row occurs once. A separate factor
checker verifies all affine pullbacks to the byte-pinned original sources.
Thus fusion, inlining and pruning are validated without trusting their
selector algorithms or cost calculations.

The final graph audit independently reconstructs square differences, every
occurrence and factor binding, every block boundary, each internal tree,
all aliases and units, and the symbolic lift divergence. Exact expansion
checks the entire height-one array, degree four, constant one and coordinate
participation. The literal text form is independently parsed as well.
Normal and optimized-Python replays, deliberate corruption tests and exact
same-source/hub/seed controls bind these properties to the delivered bytes.
A newly generated parameter-matched control is explicitly labelled new; it
is not called a historical endpoint. Full normalization comparisons are not
misrepresented as isolated one-operation experiments.

Combining the inverse maps with~\eqref{eq:networklift} and the named original
source calibration proves~\eqref{eq:calibration} for each selected polynomial.
Successful coefficient checks are not substituted for that unbounded argument.
